\documentclass[11pt]{book}
\usepackage{fontspec}
\setmainfont{Latin Modern Roman}
\setsansfont{Latin Modern Sans}
\usepackage{amsmath,amssymb,mathtools}
\usepackage{unicode-math}
\setmathfont{Latin Modern Math}
\usepackage{xurl}
\usepackage[hidelinks]{hyperref}
\newcommand{\sref}[2]{\hyperlink{source:#1:#2}{[S#2]}}
\newcommand{\cataloguescope}[1]{\paragraph{Mathematical question.}#1}
\begin{document}
\section{Research Problems in Function Theory — Problem 3.21}\label{statement}

\subsection{Definitions and mathematical statement}

\paragraph{Shared notation.}$\mathbb N=\{1,2,\ldots\}$, and $\mathbb Z,\mathbb Q,\mathbb R,\mathbb C$ denote the integers, rationals, reals and complex numbers. Any other conventions are specified locally.


Write $\mathbb D=\{z\in\mathbb C:|z|<1\}$ and $\mathbb T=\partial\mathbb D$. A \emph{continuum} is a nonempty compact connected set. For a continuum $K\subset\overline{\mathbb D}$, put $d=\operatorname{diam}K=\sup\{|z-w|:z,w\in K\}\in[0,2]$. The harmonic measure of $K$ seen from a point $z$ in the disc means the probability that planar Brownian motion started at $z$ first hits $K$ before exiting the disc, counting an exit through $K\cap\mathbb T$ as a hit. Equivalently, for $z\notin K$ it is the harmonic measure of $K$ on the boundary of the component of $\mathbb D\setminus K$ containing $z$. If $z\in K$, set this measure equal to $1$. Denote it by $\omega(z;\mathbb D;K)$.
\paragraph{Statement recorded in the supplied source.}
Is it always true that
\[\omega(0;\mathbb D;K)\ge\frac1\pi\arcsin\!\left(\frac d2\right)?\]
Here $\arcsin$ takes values in $[0,\pi/2]$ on $[0,1]$. \sref{statement}{2}
\subsection{Short English statement}
Is the probability of hitting any continuum in the closed unit disc, starting at its centre, at least the stated function of the continuum’s diameter?
\subsection{Sources}
\begin{description}
\item[\hypertarget{source:statement:1}{S1}] ULAM.AI and the UnsolvedMath Contributors, \emph{UnsolvedMath: A Curated Collection of Open Mathematics Problems}, catalogue entry. CC BY 4.0. \url{https://huggingface.co/datasets/ulamai/UnsolvedMath}.
\item[\hypertarget{source:statement:2}{S2}] Walter K. Hayman and Edward F. Lingham, \emph{Research Problems in Function Theory}, arXiv:1809.07200, 2018. Problem 3.21, its complete Update 3.21, and the relevant chapter notation. \url{https://arxiv.org/pdf/1809.07200}.
\end{description}
\label{end:statement}
\end{document}
